\documentclass[aspectratio=169, 10pt]{beamer}

\usepackage[utf8]{inputenc}
\usepackage[english]{babel}
\usepackage{amsmath, amssymb}
\usepackage{graphicx}
\usepackage{booktabs}
\usepackage{xcolor}

% --- FLORENCE UNIVERSITY COLORS ---
\definecolor{UniFiBlue}{RGB}{0, 51, 102}
\definecolor{UniFiLightBlue}{RGB}{204, 220, 235}
\definecolor{UniFiOrange}{RGB}{237, 107, 6}

% --- BEAMER THEME SETUP ---
\usetheme{Madrid}
\usecolortheme[named=UniFiBlue]{structure}

\setbeamertemplate{navigation symbols}{}
\setbeamercolor{palette primary}{bg=UniFiBlue,fg=white}
\setbeamercolor{palette secondary}{bg=UniFiBlue!80!black,fg=white}
\setbeamercolor{palette tertiary}{bg=UniFiBlue!60!black,fg=white}

\setbeamercolor{title}{bg=UniFiBlue, fg=white} 
\setbeamercolor{frametitle}{bg=UniFiBlue, fg=white}
\setbeamercolor{item projected}{bg=UniFiBlue, fg=white}

\usefonttheme{professionalfonts}

% --- TITLE PAGE INFORMATION ---
\title[Deep Learning Dynamics]{\textbf{Empirical Analysis of Learning Dynamics:}\\ Double Descent, Neural Collapse and CKA}
\author[Marco Bardazzi]{Marco Bardazzi}
\institute[UniFi]{
    \textbf{University of Florence}\\
    Computer Vision and Intelligent Media Recognition\\
    \vspace{0.3cm}
}
\date{\today}

\begin{document}

% ---------------------------------------------------
\begin{frame}
    \titlepage
\end{frame}

% ---------------------------------------------------
\begin{frame}{Presentation Agenda}
    \tableofcontents
\end{frame}

% ===================================================
\section{Theoretical Introduction}
% ===================================================

% ---------------------------------------------------
\begin{frame}{The Modern Paradigm: Deep Double Descent}
    \begin{columns}
        \begin{column}{0.5\textwidth}
            \textbf{Classical Theory:}
            \begin{itemize}
                \item ``U''-shaped Bias-Variance Trade-off.
                \item Complex models memorize noise (``overfitting'').
            \end{itemize}
            \vspace{0.3cm}
            \textbf{The Reality of Deep Learning:}
            \begin{itemize}
                \item Past the \textbf{interpolation threshold}, the test error starts decreasing again.
                \item Massively over-parameterized models generalize excellently.
            \end{itemize}
        \end{column}
        \begin{column}{0.5\textwidth}
            \begin{center}
                % --- INSERT IMAGE HERE ---
                % Crop from original file: Row 1, Column 1 (Exp 1: Test Error)
                \includegraphics[width=\linewidth]{1.png}
                %\fcolorbox{gray}{lightgray}{\parbox[c][5cm][c]{0.9\linewidth}{\centering \textbf{INSERT HERE:\\ Exp 1: Test Error\\ (Row 1, Col 1)}}}
            \end{center}
        \end{column}
    \end{columns}
\end{frame}

% ---------------------------------------------------
\begin{frame}{What is Neural Collapse (NC)?}
    In the terminal phases of training, last-layer representations exhibit 4 geometric properties (Papyan et al., 2020):
    \vspace{0.3cm}
    \begin{enumerate}
        \item \textbf{NC1 (Variability Collapse):} Activations converge to their class center (within-class variance $\to 0$).
        \item \textbf{NC2 (Simplex ETF):} Class centers arrange symmetrically, maximizing distance (Equiangular Tight Frame).
        \item \textbf{NC3 (Self-Duality):} The linear classifier weights align perfectly with the class centers.
        \item \textbf{NC4 (NCC):} Network predictions coincide with a Nearest Class Center classification.
    \end{enumerate}
\end{frame}

% ===================================================
\section{Experimental Setup}
% ===================================================

% ---------------------------------------------------
\begin{frame}{Experimental Setup}
    \textbf{Objective:} Empirically unify Double Descent, Neural Collapse, and layer similarity analysis (CKA).
    
    \vspace{0.3cm}
    \begin{itemize}
        \item \textbf{Dataset:} MNIST-1D. Low-dimensional but morphologically complex data.
        \item \textbf{Injected noise:} \textcolor{red}{\textbf{20\% Label Noise}}. Forces the network to memorize incorrect labels, exacerbating the interpolation peak.
        \item \textbf{Architectures:} \texttt{ResNet1D} and \texttt{StandardCNN1D}.
        \item \textbf{Scaling:} Variable width from $W=2$ to $W=64$ ($10^2 \to 10^5$ params).
        \item \textbf{Optimizers:} SGD vs Adam.
    \end{itemize}
\end{frame}

% ===================================================
\section{Results Analysis (Exp 1, 2, 3)}
% ===================================================

% ---------------------------------------------------
\begin{frame}{Results: The Breakdown of Collapse in the Critical Regime}
    \begin{columns}
        \begin{column}{0.5\textwidth}
            \begin{itemize}
                \item Generalization error correlates perfectly with the \textbf{NC1} metric.
                \item Around $10^3-10^4$ parameters, within-class variance explodes. The model breaks its geometry to memorize noise.
                \item Past the threshold ($>10^4$), Neural Collapse is restored and the error drops.
            \end{itemize}
        \end{column}
        \begin{column}{0.5\textwidth}
            \begin{center}
                % --- INSERT IMAGE HERE ---
                % Crop from original file: Row 1, Column 2 (Exp 1: NC1)
                 \includegraphics[width=\linewidth]{2.png}
                %\fcolorbox{gray}{lightgray}{\parbox[c][5cm][c]{0.9\linewidth}{\centering \textbf{INSERT HERE:\\ Exp 1: NC1 (Variability)\\ (Row 1, Col 2)}}}
            \end{center}
        \end{column}
    \end{columns}
\end{frame}

% ---------------------------------------------------
\begin{frame}{The Impact of the Optimizer: SGD dominates Adam}
    \begin{columns}
        \begin{column}{0.5\textwidth}
            \begin{itemize}
                \item In the ``overparameterized'' regime, \textbf{SGD} achieves much better geometric configurations (NC2).
                \item \textbf{Adam} tends to get stuck in geometrically sub-optimal local minima.
                \item \textbf{Why?} SGD possesses an \textit{implicit inductive bias} that penalizes weight norms and favors wide, symmetric margins.
            \end{itemize}
        \end{column}
        \begin{column}{0.5\textwidth}
            \begin{center}
                % --- INSERT IMAGE HERE ---
                % Crop from original file: Row 1, Column 3 (Exp 1: NC2)
                 \includegraphics[width=\linewidth]{3.png}
                %\fcolorbox{gray}{lightgray}{\parbox[c][5cm][c]{0.9\linewidth}{\centering \textbf{INSERT HERE:\\ Exp 1: NC2 (Simplex ETF)\\ (Row 1, Col 3)}}}
            \end{center}
        \end{column}
    \end{columns}
\end{frame}

% ---------------------------------------------------
\begin{frame}{Training Dynamics (Epoch-wise)}
    \begin{columns}
        \begin{column}{0.5\textwidth}
            Comparison between two configurations:
            \vspace{0.2cm}
            \begin{itemize}
                \item \textbf{Critical ($W=12$):} Located on the Double Descent peak. The loss surface is highly fragmented, leading to massive error oscillations.
                \item \textbf{Large ($W=64$):} Excess capacity allows rapid learning of robust features, relegating noise memorization without destroying generalization.
            \end{itemize}
        \end{column}
        \begin{column}{0.5\textwidth}
            \begin{center}
                % --- INSERT IMAGE HERE ---
                % Crop from original file: Row 2, Column 2 (Exp 2: Epoch-wise Test Error)
                 \includegraphics[width=\linewidth]{4.png}
                %\fcolorbox{gray}{lightgray}{\parbox[c][5cm][c]{0.9\linewidth}{\centering \textbf{INSERT HERE:\\ Exp 2: Epoch-wise Error\\ (Row 2, Col 2)}}}
            \end{center}
        \end{column}
    \end{columns}
\end{frame}

% ---------------------------------------------------
\begin{frame}{Dimensional Limits of Neural Collapse}
    \begin{columns}
        \begin{column}{0.5\textwidth}
            \begin{itemize}
                \item Mathematically, a perfect Simplex ETF with $K$ classes requires a latent dimension of \textbf{$d \ge K-1$}.
                \item For MNIST-1D ($K=10$), we expect collapse starting from $d=9$.
                \item \textbf{Empirical Verification:} NC2 drops vertically and stabilizes at zero exactly at the theoretical limit.
            \end{itemize}
        \end{column}
        \begin{column}{0.5\textwidth}
            \begin{center}
                % --- INSERT IMAGE HERE ---
                % Crop from original file: Row 2, Column 4 (Exp 3: Gen NC NC2 vs dim)
                 \includegraphics[width=\linewidth]{5.png}
                %\fcolorbox{gray}{lightgray}{\parbox[c][5cm][c]{0.9\linewidth}{\centering \textbf{INSERT HERE:\\ Exp 3: NC2 vs dim\\ (Row 2, Col 4)}}}
            \end{center}
        \end{column}
    \end{columns}
\end{frame}

% ===================================================
\section{Representations Analysis (CKA)}
% ===================================================

% ---------------------------------------------------
\begin{frame}{CKA: Structural Differences (Small vs Large)}
    \begin{columns}
        \begin{column}{0.5\textwidth}
            \textbf{CKA:} Measures geometric similarity ($0 \to 1$).
            \vspace{0.2cm}
            \begin{itemize}
                \item \textbf{L0/L1:} High CKA. Critical and over-parameterized networks extract the same low-level features.
                \item \textbf{L3 (Deep):} CKA drops ($\approx 0.34$).
                \item The small model lacks the degrees of freedom to create the Simplex ETF, resulting in a diametrically opposed final topology.
            \end{itemize}
        \end{column}
        \begin{column}{0.5\textwidth}
            \begin{center}
                % --- INSERT IMAGE HERE ---
                % Crop from original file: Row 3, Column 1 (Exp 4A: Heatmap CKA)
                 \includegraphics[width=\linewidth]{6.png}
                %\fcolorbox{gray}{lightgray}{\parbox[c][5cm][c]{0.9\linewidth}{\centering \textbf{INSERT HERE:\\ Exp 4A: Heatmap\\ (Row 3, Col 1)}}}
            \end{center}
        \end{column}
    \end{columns}
\end{frame}

% ---------------------------------------------------
\begin{frame}{CKA: Drift Over Time (Epoch Drift)}
    \begin{columns}
        \begin{column}{0.5\textwidth}
            Comparison of activations in successive epochs relative to epoch 1:
            \vspace{0.2cm}
            \begin{itemize}
                \item \textbf{Stem (L0):} Converges immediately (CKA constant near 1).
                \item \textbf{Deep Layers (L2, L3):} Continuous drift. The latent space reorganizes over thousands of epochs, driven by the minimization of within-class variance.
            \end{itemize}
        \end{column}
        \begin{column}{0.5\textwidth}
            \begin{center}
                % --- INSERT IMAGE HERE ---
                % Crop from original file: Row 3, Column 2 (Exp 4B: CKA Line chart)
                 \includegraphics[width=\linewidth]{7.png}
                %\fcolorbox{gray}{lightgray}{\parbox[c][5cm][c]{0.9\linewidth}{\centering \textbf{INSERT HERE:\\ Exp 4B: Epoch-wise\\ (Row 3, Col 2)}}}
            \end{center}
        \end{column}
    \end{columns}
\end{frame}

% ---------------------------------------------------
\begin{frame}{CKA: Reproducibility and Geometric Degeneracy}
    \begin{columns}
        \begin{column}{0.5\textwidth}
            Same model ($W=32$), only \textit{different initial random Seeds}:
            \vspace{0.2cm}
            \begin{itemize}
                \item L0 learns deterministic features ($CKA = 0.97$).
                \item \textbf{L3 varies significantly ($CKA = 0.59$).}
                \item \textbf{Geometric Degeneracy:} There are infinite orthogonal rotations to create a perfect Simplex. Networks converge to logically isomorphic solutions, but in different absolute coordinates.
            \end{itemize}
        \end{column}
        \begin{column}{0.5\textwidth}
            \begin{center}
                % --- INSERT IMAGE HERE ---
                % Crop from original file: Row 3, Column 3 (Exp 4C: Bar chart)
                 \includegraphics[width=\linewidth]{8.png}
                %\fcolorbox{gray}{lightgray}{\parbox[c][5cm][c]{0.9\linewidth}{\centering \textbf{INSERT HERE:\\ Exp 4C: Cross-Seed\\ (Row 3, Col 3)}}}
            \end{center}
        \end{column}
    \end{columns}
\end{frame}

% ===================================================
\section{Conclusions}
% ===================================================

% ---------------------------------------------------
\begin{frame}{Conclusions}
    \begin{enumerate}
        \item \textbf{The Empirical Link:} The Double Descent phenomenon is intimately linked to the breakdown (in the critical regime) and subsequent reconstruction (in the massive regime) of Neural Collapse.
        \item \textbf{Optimization:} SGD dominates Adam, acting as an implicit regularizer that guides networks toward a perfect Simplex ETF.
        \item \textbf{Deep Geometry:} CKA revealed the heterogeneity of learning. While early layers are static and deterministic, the last layer degenerates geometrically, undergoing a slow reorganization to maximize margins.
    \end{enumerate}
\end{frame}

% ---------------------------------------------------
\begin{frame}
    \begin{center}
        \Huge \textbf{Thank you for your attention!}\\
        \vspace{1cm}
        \Large Questions?
    \end{center}
\end{frame}

\end{document}